% CARE Conference 2026, IIT Roorkee -- 1-page extended abstract
% Poster track submission. Compile: pdflatex CARE_AI2026_Abstract.tex
\documentclass[10pt,a4paper]{article}
\usepackage[margin=0.7in]{geometry}
\usepackage[T1]{fontenc}
\usepackage{times}
\usepackage[hidelinks]{hyperref}
\usepackage{microtype}
\usepackage{parskip}
\pagestyle{empty}
\setlength{\parskip}{0.32em}
\linespread{0.97}

\newcommand{\sect}[1]{\vspace{0.25em}\noindent\textbf{#1}\par\nobreak\vspace{-0.15em}}

\begin{document}

\begin{center}
{\large\bfseries Label Reliability, Not Model Capacity, Limits Longitudinal\\
Prediction of Parkinson's Disease Motor Subtypes}\\[0.6em]
Parv Patel, Mayank Tewatia, and Dr.\ Sahasranand K.~R.\\[0.2em]
\emph{Indian Institute of Technology Palakkad, India}
\end{center}
\vspace{0.3em}

\sect{1. Motivation}
Parkinson's disease is commonly divided into tremor-dominant (TD) and postural
instability and gait difficulty (PIGD) motor subtypes, with an Indeterminate band
between them. The subtype is obtained by thresholding the ratio of two averages of
MDS-UPDRS items [1], and a change of subtype over time is widely treated as a marker
of motor progression. Much modelling work therefore aims to predict this label, but
the reliability of the label itself is rarely measured first. This study asks whether
longitudinal prediction of motor subtype is limited by the choice of model family, or
by the label being predicted.

\sect{2. Data and evaluation protocol}
We used the Parkinson's Progression Markers Initiative (PPMI) observational cohort:
1,080 patients and 4,441 labelled visits, of which 934 patients with two or more
labelled visits contribute 3,361 consecutive visit-to-visit transitions.

We applied one protocol to four model families: (i) discrete-time and continuous-time
Markov chains on the observed label; (ii) a hidden Markov model (HMM) with a latent
state and a fitted misclassification matrix; (iii) a hierarchical Bayesian growth
model on the ratio; and (iv) Gaussian process (GP) regression on the ratio and on the
tremor and PIGD subscores.

Every model was evaluated by leave-one-patient-out validation. No visit of the
held-out patient enters the fit, and each prediction uses only that patient's own
earlier visits. All intervals are patient-clustered bootstrap intervals, since
predictions are correlated within patients. A ten-check battery verified fold
construction, history ordering and prediction counts on every saved prediction row.

\sect{3. Results}
\emph{Label stability.} 35.2\% of all observed subtype changes reverted at the very
next visit. Requiring two consecutive assessments before accepting a conversion
removed 51.7\% of apparent PIGD conversions (263 to 127). The Markov assumption on the
observed label was rejected in all three current states ($p \le 0.004$), which
motivates a latent-state formulation.

\emph{Discrete-state models.} Held-out accuracy over the same 3,361 transitions: HMM
0.7539 [0.7348, 0.7731]; discrete Markov chain 0.7474 [0.7272, 0.7670]; continuous-time
chain 0.7453 [0.7250, 0.7646]; persistence, meaning no change, 0.7227 [0.7009, 0.7437];
marginal null 0.6293 [0.6011, 0.6565]. The HMM is the best model, but its paired gain
over the rule that simply repeats the current subtype is only $+3.1$ points (95\% CI
$+2.2$ to $+4.1$), and it assigns the Indeterminate class once in 3,361 opportunities.
Increasing the latent state space from three to five did not improve accuracy (0.7513).

\emph{Why the ceiling is where it is.} The hierarchical Bayesian model estimates a
population trend of $-0.093$ per year in the ratio (95\% credible interval $-0.116$ to
$-0.068$) but a total residual spread of 1.14, against an Indeterminate band only 0.25
wide, a factor of about 4.6. The middle class is therefore almost unreachable for any
model predicting the band from a continuous quantity.

\emph{Continuous-trajectory models.} On the ratio, a GP with a flexible kernel improves
on a linear-trend GP ($-0.033$ MAE, 95\% CI $-0.047$ to $-0.021$) and on persistence.
But when the two subscores were modelled separately, the GP did not beat persistence on
tremor ($-0.0031$, 95\% CI $-0.0084$ to $+0.0022$) and was worse on PIGD ($+0.0134$,
95\% CI $+0.0083$ to $+0.0187$).

\emph{Matched comparison.} Converting the subscore GP predictions back to a subtype
allows a like-for-like comparison with the HMM on exactly the same 934 patients and
3,361 points. GP-derived accuracy was 75.66\% against the HMM's 75.39\%, a paired
difference of $+0.27$ points (95\% CI $-0.42$ to $+0.98$).

\sect{4. Discussion and limitations}
Four model families with substantially different assumptions converge on the same
held-out performance, only a few points above a trivial rule, and none recovers the
Indeterminate class. With frequent label reversion and a residual spread several times
the band width, this indicates that the representation of motor symptoms, rather than
model capacity, is the binding constraint.

The target is our own rule's next recorded label, not an independent clinical
diagnosis. The ratio is undefined when the PIGD subscore is zero, excluding 692 visits
that are overwhelmingly TD, so the ratio-based analyses are not a random subset. The
matched comparison is a failure to detect a difference and does not establish
equivalence. No external cohort was used and no clinical utility is claimed. Ongoing
work models the two subscores jointly under a bounded likelihood, which addresses both
the undefined-ratio exclusion and predictive mass below zero.

{\footnotesize
\vspace{0.3em}
\noindent\textbf{Reference.} [1] Stebbins GT, Goetz CG, Burn DJ, Jankovic J, Khoo TK,
Tilley BC. How to identify tremor dominant and postural instability/gait difficulty
groups with the Movement Disorder Society Unified Parkinson's Disease Rating Scale.
\emph{Movement Disorders} 2013;28(5):668--70.

\vspace{0.2em}
\noindent\textbf{Data.} Parkinson's Progression Markers Initiative (PPMI),
\url{ppmi-info.org}. % TODO: insert the current PPMI acknowledgement wording here.
}

\end{document}
